% Source: https://artofproblemsolving.com/community/c3h1601880_number_problems
% Stable identity: aops-topic-1601880-q02
Let $77n=k(k+1)(k+2)$ with $k\ge1$. Since $7$ and $11$ are primes, one of the three consecutive numbers must be a multiple of $7$ and one must be a multiple of $11$ (possibly the same one, if it is a multiple of $77$). As $k$ increases the product increases, so we want the smallest such $k$.

The multiples of $7$ are $7,14,21,\dots$ and of $11$ are $11,22,\dots$. A multiple of $7$ and a multiple of $11$ within distance $2$ of each other first occur at $21$ and $22$ ($7$ and $11$ differ by $4$, $14$ and $11$ by $3$); the first multiple of $77$ is $77$ itself, which is larger. So the smallest $k$ is $k=20$, giving
\[
20\cdot21\cdot22=9240=77\cdot120 .
\]
Hence $n=\boxed{120}$.
